\scanpage{6}
% archive sheet number: 5

on trouve un $\underline{\unsure{I|J|\mathcal J}}_{Y_0}$ qu. coh. sur $Y_0$.

Pour prouver la M-effectivité de $R$, on peut supposer que $Y_0$ \deleted{\unsure{devient|descente} \unsure{$Y$|$X$}} affine. Soit alors $A_0$ l'anneau de $Y_0$, et \inserted{soit} \deleted{$A=\check H^0(X/\unsure{(X/R)|R}),\underline{\mathcal O}_X)$}

\[
A=\operatorname{Ker}\left(H^0(X,\mathcal O_X)\rightrightarrows H^0(R,\underline{\mathcal O}_R)\right)
\]

On a un hom évident $A\longrightarrow A_0$, dont le noyau est $\operatorname{Ker}\bigl(H^0(X_0,\underline I_{X_0})\rightrightarrows H^0(R_0,I_{R_0})\bigr)=H^0(Y_0,\underline I_{Y_0})=I_{Y_0}$. \deleted{\unsure{Il y a une suite}} \inserted{Je dis que la} \deleted{\unsure{suite ex}} suite d'hom

\[
0\longrightarrow I_{Y_0}\longrightarrow A\longrightarrow A_0\longrightarrow0
\]

est \uline{exacte}, i.e. $A\longrightarrow A_0$ \uline{surjectif}.

Pour le voir, \deleted{prenons} soit $\varphi_0\in A_0$, soit $f_0$ son image \inserted{inverse} dans $X_0$, \deleted{$f_0\in H^0(X_0,\mathcal O_{X_0})$} \inserted{donc}

\[
f_0\in\operatorname{Ker}\left(H^0(X_0,\mathcal O_{X_0})\rightrightarrows H^0(R_0,\mathcal O_{R_0})\right)
\]

Je dis d'abord que $f_0$ se remonte en un $f\in H^0(X,\underline{\mathcal O}_X)$. En effet, à priori il y a une obstruction à relever qui est \deleted{dans} $\unsure{\zeta|\xi}\in H^1(X_0,I_{X_0})$ ; mais les deux images
